\documentclass{article}

% Core packages
\usepackage[utf8]{inputenc}
\usepackage[T1]{fontenc}
\usepackage{amsmath,amssymb,amsfonts}
\usepackage{mathtools}
\usepackage{bm}
\usepackage{booktabs}
\usepackage{multirow}
\usepackage{graphicx}
\usepackage{xcolor}
\usepackage{hyperref}
\usepackage{cleveref}
\usepackage{algorithm}
\usepackage{algorithmic}
\usepackage{subcaption}
\usepackage{enumitem}
\usepackage[margin=1in]{geometry}
\usepackage{natbib}
\usepackage{amsthm}

\newtheorem{proposition}{Proposition}

% Custom commands
\newcommand{\vx}{\bm{x}}
\newcommand{\vo}{\bm{o}}
\newcommand{\vz}{\bm{z}}
\newcommand{\va}{\bm{a}}
\newcommand{\vep}{\bm{\epsilon}}
\newcommand{\vh}{\bm{h}}
\newcommand{\R}{\mathbb{R}}
\newcommand{\E}{\mathbb{E}}
\newcommand{\KL}{\mathrm{KL}}
\newcommand{\loss}{\mathcal{L}}
\DeclareMathOperator{\sigmoid}{sigmoid}
\DeclareMathOperator{\diag}{diag}

\title{Action-Conditioned Unified Latents for Robotic World Models}

\author{
  Anonymous Authors
}

\date{}

\begin{document}

\maketitle

\begin{abstract}
Learning compact, physics-grounded visual representations is a central challenge for building robotic world models that jointly predict future observations and actions.
Existing approaches either reconstruct raw pixels---wasting capacity on task-irrelevant details---or use frozen pretrained features that discard action-relevant geometric information.
We introduce \textbf{Action-Conditioned Unified Latents (AC-UL)}, a framework that extends the Unified Latents (UL) variational objective to autoregressive state-action dynamics.
Given a history of observations $\vo_{t-H+1:t}$ and past actions $\va_{t-H:t-1}$, our dynamics prior \emph{jointly} generates the next observation latent $\vz_{t+1}$ and the current action $\va_t$---unifying world modeling and policy learning in a single generative process.
By conditioning the diffusion prior on the observation-action history, the variational bound \emph{mathematically forces} the encoder to discard unpredictable visual content while retaining the physical geometry that the history of interactions can predict.
In Stage~1, the encoder, dynamics prior, and diffusion decoder are jointly trained; in Stage~2, the frozen encoder provides a physics-grounded visual vocabulary on which an autoregressive diffusion transformer is trained for joint video-action generation, enabling closed-loop planning via iterative rollouts.
We evaluate AC-UL on simulated manipulation benchmarks (CALVIN, LIBERO) and real-world robotic video (DROID), demonstrating improvements over pixel-reconstruction world models and frozen-feature approaches in both planning accuracy and policy performance.
\end{abstract}

%=====================================================================
\section{Introduction}
\label{sec:intro}
%=====================================================================

Building world models---internal simulators that predict how actions change the environment---is a long-standing goal in robotics and reinforcement learning~\citep{ha2018world,hafner2023dreamerv3}.
A world model that operates in a learned latent space can be orders of magnitude cheaper to roll out than a pixel-level simulator, enabling model-predictive control, planning, and data-efficient policy learning.

Recent progress in autoregressive video-action models~\citep{dreamzero2025} has shown that jointly generating future video frames and robot actions from a shared backbone can achieve remarkable generalization.
Meanwhile, latent diffusion models~\citep{rombach2022high} and vision-language-action (VLA) architectures~\citep{brohan2023rt2,kim2024openvla} have produced powerful visual representations.
However, two fundamental problems remain:

\textbf{(1) Pixel-reconstruction world models} (e.g., DreamerV3~\citep{hafner2023dreamerv3}, DreamZero~\citep{dreamzero2025}) spend capacity encoding textures, lighting, and background clutter that are irrelevant to the dynamics.
The learned latent confounds physical geometry with appearance, making long-horizon rollouts noisy and planning unreliable.

\textbf{(2) Frozen-feature approaches} (e.g., DINO-WM~\citep{zhou2024dinowm}) use pretrained encoders (DINOv2) that were not trained to understand action effects.
While these features transfer surprisingly well, they are agnostic to the causal structure of manipulation.

We propose a third path: \textbf{learning the latent representation itself to encode only the information that the history of interactions can predict}.
Our approach extends the Unified Latents (UL) framework~\citep{heek2025ul} from unconditional image modeling to \emph{joint} observation-action dynamics.
The key insight is that by conditioning the diffusion prior on the history of past observations and actions, and requiring it to jointly predict the next observation latent \emph{and} the current action, any encoder information that cannot be predicted from this history will spike the variational loss.
The encoder is thus mathematically forced to discard unpredictable pixels and retain only the physical geometry that the robot's interaction history can explain.

This formulation---where the model takes in $(\vo_{t-H+1:t}, \va_{t-H:t-1})$ and generates $(\vo_{t+1}, \va_t)$---follows the autoregressive structure of DreamZero~\citep{dreamzero2025}, but operates in a principled latent space rather than raw pixels.

Our contributions are:
\begin{enumerate}[leftmargin=*,itemsep=2pt]
  \item We derive the \textbf{Action-Conditioned Unified Latent (AC-UL)} objective, extending the UL variational bound to history-conditioned \emph{joint} observation-action dynamics (\Cref{sec:method}).
  \item We provide a complete \textbf{two-stage training algorithm}: Stage~1 jointly trains the encoder, history-conditioned dynamics prior, and diffusion decoder; Stage~2 trains an autoregressive diffusion transformer on frozen latents for joint video-action generation (\Cref{sec:stage2}).
  \item We propose and evaluate three downstream applications: \textbf{closed-loop planning via latent rollouts}, \textbf{imitation learning with joint world-action modeling} (compared to FLARE~\citep{nvidia2025flare}), and \textbf{real-world video prediction} on the DROID dataset (\Cref{sec:experiments}).
\end{enumerate}

%=====================================================================
\section{Background}
\label{sec:background}
%=====================================================================

\subsection{Variational Autoencoders and the ELBO}

Given data $\vx$ and latent variable $\vz_0$, the Evidence Lower Bound (ELBO) on the log-likelihood is:
\begin{equation}
  -\log p_\theta(\vx) \leq \E_{\vz_0 \sim p_\theta(\vz_0|\vx)}\!\Big[\underbrace{-\log p_\theta(\vx|\vz_0)}_{\text{decoder}}\Big] + \KL\!\Big[\underbrace{p_\theta(\vz_0|\vx)}_{\text{encoder}} \,\Big\|\, \underbrace{p_\theta(\vz_0)}_{\text{prior}}\Big] = \loss(\vx).
  \label{eq:elbo}
\end{equation}
In the Unified Latents framework, both the decoder $p_\theta(\vx|\vz_0)$ and the prior $p_\theta(\vz_0)$ are parameterized as diffusion models.

\subsection{Diffusion Models and the KL Bound}

Consider a destruction process $\vx_t = \alpha(t)\,\vx + \sigma(t)\,\vep$ with $\vep \sim \mathcal{N}(\bm{0},\bm{I})$, where the noise level is characterized by the log signal-to-noise ratio $\lambda(t) = \log(\alpha_t^2/\sigma_t^2)$ with the variance-preserving constraint $\alpha_t^2 + \sigma_t^2 = 1$.
The KL divergence between the data-conditioned diffusion path and the learned reverse process can be bounded as~\citep{ho2020denoising,kingma2021variational}:
\begin{equation}
  \KL\!\Big[p(\vx_0|\vx) \,\Big\|\, p_\theta(\vx_0)\Big]
  \leq \E_{t\sim\mathcal{U}(0,1)}\!\left[-\frac{d\lambda(t)}{dt}\,\frac{\exp\lambda(t)}{2}\,w(\lambda_t)\,\|\vx - \hat{\vx}(\vx_t,\theta)\|^2\right]
  + \KL\!\Big[p(\vx_1|\vx) \,\Big\|\, \mathcal{N}(\bm{0},\bm{I})\Big],
  \label{eq:diff_kl}
\end{equation}
where $w(\lambda_t) = 1$ yields a tight bound, and $w(\lambda_t) = \sigmoid(\lambda_t - b)$ provides a perceptually motivated re-weighting.

\subsection{Unified Latents (UL)}
\label{sec:ul_background}

UL~\citep{heek2025ul} uses a deterministic encoder $\vz_{\text{clean}} = E_\theta(\vx)$ with fixed Gaussian noise added up to a minimum noise level defined by $\lambda_z(0) = 5$:
\begin{equation}
  \vz_0 = \alpha_z(0)\,\vz_{\text{clean}} + \sigma_z(0)\,\vep_z, \quad
  \alpha_z(0) = \sqrt{\sigmoid(+5)} \approx 1.0, \quad
  \sigma_z(0) = \sqrt{\sigmoid(-5)} \approx 0.08.
  \label{eq:ul_noise}
\end{equation}

\textbf{Prior loss (Equation~3 of UL).}
The KL penalty between the encoder distribution and the diffusion prior is:
\begin{equation}
  \KL\!\Big[p(\vz_0|\vx) \,\Big\|\, p_\theta(\vz_0)\Big]
  \leq \E_{t}\!\left[-\frac{d\lambda_z(t)}{dt}\,\frac{\exp\lambda_z(t)}{2}\,w(\lambda_z(t))\,\|\vz_{\text{clean}} - \hat{\vz}(\vz_t, \theta)\|^2\right]
  + \KL\!\Big[p(\vz_1|\vx) \,\Big\|\, \mathcal{N}(\bm{0},\bm{I})\Big],
  \label{eq:ul_prior}
\end{equation}
where $w(\lambda_z(t)) = 1$ for a tight bitrate bound.

\textbf{Decoder loss (Equation~4 of UL).}
The reconstruction term uses a diffusion decoder conditioned on $\vz_0$:
\begin{equation}
  -\log p_\theta(\vx|\vz_0) \leq \E_{t\sim\mathcal{U}(0,1)}\!\left[\frac{d\lambda_x(t)}{dt}\,\frac{\exp\lambda_x(t)}{2}\,w_x(\lambda_x(t))\,\|\vx - \hat{\vx}(\vx_t, \vz_0, \theta)\|^2\right],
  \label{eq:ul_decoder}
\end{equation}
with sigmoid weighting $w_x(\lambda_x(t)) = \sigmoid(\lambda_x(t) - b)$ and a loss factor $c_{\text{lf}} \in [1.3, 1.7]$ that up-weighs the decoder to prevent posterior collapse.

\subsection{Autoregressive Video-Action Models}
\label{sec:dreamzero_background}

DreamZero~\citep{dreamzero2025} introduced an autoregressive diffusion transformer that jointly generates video frames and robot actions.
Given a history of observations $\vo_{t-H+1:t}$ and past actions $\va_{t-H:t-1}$, the model generates the next observation $\vo_{t+1}$ and the current action $\va_t$ as a single denoising step:
\begin{equation}
  p_\theta(\vo_{t+1}, \va_t \mid \vo_{t-H+1:t}, \va_{t-H:t-1}).
  \label{eq:dreamzero}
\end{equation}
This can be decomposed as video prediction followed by inverse dynamics: $p_\theta(\vo_{t+1} \mid \text{history}) \times p_\theta(\va_t \mid \vo_{t+1}, \text{history})$.
During closed-loop execution, predicted frames are replaced with ground-truth observations after each action, preventing error accumulation.
We adopt this autoregressive structure but operate in a learned latent space that is optimized for action-relevant content.

%=====================================================================
\section{Method: Action-Conditioned Unified Latents}
\label{sec:method}
%=====================================================================

We extend the UL framework from single-image encoding to \emph{history-conditioned joint observation-action dynamics}.
The central idea is to replace the unconditional prior $p_\theta(\vz_0)$ with a \emph{dynamics prior} $p_\theta(\vz_{t+1,0}, \va_t \mid \mathcal{H}_t)$ that must jointly predict the future observation latent and the current action from the interaction history.

\subsection{Problem Setup and Variables}
\label{sec:variables}

We consider a robotic agent that observes image frames and executes actions in an autoregressive loop. We define the \emph{interaction history} at time $t$:
\begin{equation}
  \mathcal{H}_t = \big(\vo_{t-H+1}, \va_{t-H}, \vo_{t-H+2}, \va_{t-H+1}, \ldots, \vo_{t-1}, \va_{t-1}, \vo_t\big),
  \label{eq:history}
\end{equation}
where:
\begin{itemize}[leftmargin=*,itemsep=2pt]
  \item $\vo_k \in \R^{H_{\text{img}} \times W_{\text{img}} \times 3}$: RGB observation at time step $k$.
  \item $\va_k \in \R^{d_a}$: robot action executed at step $k$ (e.g., end-effector displacement + gripper state).
  \item $H$: history window length.
\end{itemize}

The model's task is to jointly generate:
\begin{equation}
  (\vo_{t+1}, \va_t) \sim p_\theta(\vo_{t+1}, \va_t \mid \mathcal{H}_t).
  \label{eq:joint_gen}
\end{equation}
That is, given the current and past observations and all past actions, the model simultaneously predicts \emph{what action to take now} ($\va_t$) and \emph{what will happen next} ($\vo_{t+1}$).

\subsection{Latent Encoding and the Forward Process}
\label{sec:encoding}

The encoder $E_\theta$ compresses each observation into a deterministic latent:
\begin{equation}
  \vz_{k,\text{clean}} = E_\theta(\vo_k) \in \R^{h \times w \times c}, \quad k = t-H+1, \ldots, t, t+1,
  \label{eq:encode}
\end{equation}
with spatial dimensions $h \times w$ downsampled from $H_{\text{img}} \times W_{\text{img}}$ and $c$ latent channels.

The latent history is:
\begin{equation}
  \mathcal{H}_t^z = \big(\vz_{t-H+1,\text{clean}}, \va_{t-H}, \vz_{t-H+2,\text{clean}}, \va_{t-H+1}, \ldots, \vz_{t,\text{clean}}\big).
  \label{eq:latent_history}
\end{equation}

Following UL, we add fixed Gaussian noise to the future clean latent. For any diffusion time $\tau \in [0,1]$:
\begin{equation}
  \vz_{t+1,\tau} = \alpha_z(\tau)\,\vz_{t+1,\text{clean}} + \sigma_z(\tau)\,\vep, \quad \vep \sim \mathcal{N}(\bm{0}, \bm{I}),
  \label{eq:forward_latent}
\end{equation}
where the log-SNR schedule $\lambda_z(\tau) = \log(\alpha_z^2(\tau)/\sigma_z^2(\tau))$ satisfies $\alpha_z^2(\tau) + \sigma_z^2(\tau) = 1$.

At $\tau = 0$ (minimum noise level), $\lambda_z(0) = 5$ gives:
\begin{equation}
  \alpha_z(0) = \sqrt{\sigmoid(5)} \approx 0.9967, \quad
  \sigma_z(0) = \sqrt{\sigmoid(-5)} \approx 0.0813.
  \label{eq:min_noise}
\end{equation}

\subsection{History-Conditioned Joint Dynamics Prior}
\label{sec:ac_prior}

This is the central contribution. In standard UL, the prior $\hat{\vz}(\vz_t, \theta)$ unconditionally denoises the latent.
We replace this with a \emph{joint dynamics prior} that conditions on the full interaction history $\mathcal{H}_t^z$ and simultaneously predicts the future observation latent and the current action:
\begin{equation}
  \big(\hat{\vz}_{t+1,\text{clean}},\; \hat{\va}_t\big) = \hat{f}_\theta^{\text{dyn}}\!\left(\vz_{t+1,\tau},\; \mathcal{H}_t^z\right).
  \label{eq:dynamics_prior_pred}
\end{equation}

The prior predicts:
\begin{enumerate}[leftmargin=*,itemsep=2pt]
  \item \textbf{Future latent:} $\hat{\vz}_{t+1,\text{clean}} \approx \vz_{t+1,\text{clean}}$ --- what the next observation will look like in latent space.
  \item \textbf{Current action:} $\hat{\va}_t \approx \va_t$ --- what action should be taken now.
\end{enumerate}

\subsubsection{Derivation of the History-Conditioned Joint KL Bound}

We derive the AC-UL prior loss by applying the diffusion KL bound (\Cref{eq:diff_kl}) to the conditional distributions. The KL between the encoder's output distribution and the dynamics prior over the \emph{joint} observation-action space is:
\begin{align}
  &\KL\!\Big[p(\vz_{t+1,0}, \va_t \mid \vo_{t+1}, \mathcal{H}_t) \;\Big\|\; p_\theta(\vz_{t+1,0}, \va_t \mid \mathcal{H}_t^z)\Big].
  \label{eq:joint_kl}
\end{align}

Since the action $\va_t$ is deterministic given the data (it is observed), the encoder distribution over the joint space factorizes as:
\begin{equation}
  p(\vz_{t+1,0}, \va_t \mid \vo_{t+1}, \mathcal{H}_t) = p(\vz_{t+1,0} \mid \vo_{t+1}) \cdot \delta(\va_t),
  \label{eq:encoder_factorize}
\end{equation}
where $\delta(\va_t)$ is a point mass at the observed action.

Applying the chain rule of KL and the diffusion bound to the latent component:
\begin{align}
  &\KL\!\Big[p(\vz_{t+1,0}, \va_t \mid \vo_{t+1}, \mathcal{H}_t) \;\Big\|\; p_\theta(\vz_{t+1,0}, \va_t \mid \mathcal{H}_t^z)\Big] \nonumber \\
  &= \KL\!\Big[p(\vz_{t+1,0} \mid \vo_{t+1}) \;\Big\|\; p_\theta(\vz_{t+1,0} \mid \mathcal{H}_t^z)\Big] - \E_{p(\vz_{t+1,0}|\vo_{t+1})}\!\Big[\log p_\theta(\va_t \mid \vz_{t+1,0}, \mathcal{H}_t^z)\Big] + \text{const}.
  \label{eq:joint_kl_decompose}
\end{align}

The first term (latent KL) is bounded via the diffusion decomposition:
\begin{align}
  &\KL\!\Big[p(\vz_{t+1,0} \mid \vo_{t+1}) \;\Big\|\; p_\theta(\vz_{t+1,0} \mid \mathcal{H}_t^z)\Big] \nonumber \\
  &\leq \E_{\tau \sim \mathcal{U}(0,1)}\!\left[
    -\frac{d\lambda_z(\tau)}{d\tau}\,\frac{\exp\lambda_z(\tau)}{2}\,
    \left\|\vz_{t+1,\text{clean}} - \hat{\vz}_\theta^{\text{dyn}}\!\left(\vz_{t+1,\tau},\, \mathcal{H}_t^z\right)\right\|^2
  \right] \nonumber \\
  &\quad + \KL\!\Big[p(\vz_{t+1,1} \mid \vo_{t+1}) \;\Big\|\; \mathcal{N}(\bm{0},\bm{I})\Big].
  \label{eq:ac_prior_loss}
\end{align}

The second term (action likelihood) is trained with a standard regression loss:
\begin{equation}
  \loss_a(\theta) = \left\|\va_t - \hat{\va}_\theta\!\left(\vz_{t+1,\tau}, \mathcal{H}_t^z\right)\right\|^2.
  \label{eq:action_loss}
\end{equation}

The complete prior loss combines both:
\begin{equation}
  \loss_z^{\text{joint}}(\theta) = \loss_z^{\text{latent}}(\theta) + \gamma \cdot \loss_a(\theta),
  \label{eq:joint_prior_loss}
\end{equation}
where $\gamma$ controls the relative weight of action prediction and $\loss_z^{\text{latent}}$ is given by \Cref{eq:ac_prior_loss}.

\subsubsection{Why History-Conditioned Joint Prediction Forces Physics-Grounded Encoding}

The loss is driven by two MSE terms: (i) the latent prediction error $\|\vz_{t+1,\text{clean}} - \hat{\vz}_{t+1}\|^2$ and (ii) the action prediction error $\|\va_t - \hat{\va}_t\|^2$.

Consider what happens if the encoder $E_\theta$ encodes information into $\vz_{t+1,\text{clean}}$ that is \emph{unpredictable} from the history $\mathcal{H}_t^z$:

\begin{enumerate}[leftmargin=*]
  \item \textbf{Random textures/lighting:} If $E_\theta$ preserves random background variations, the dynamics prior cannot predict these from the history (past frames and actions do not determine future lighting). The latent MSE spikes.

  \item \textbf{Novel objects entering the scene:} Sudden appearance of objects not visible in the history is unpredictable. The encoder learns to down-weight such transients.

  \item \textbf{Action-caused geometry:} Gripper position, object displacement, and contact geometry \emph{are} predictable from the history of states and actions via learned physics. Encoding this \emph{decreases} both the latent and action MSE.
\end{enumerate}

The \emph{joint} prediction of $(\vz_{t+1}, \va_t)$ creates an additional pressure: the latent must contain enough action-relevant structure that the prior can infer \emph{what action was taken} (inverse dynamics). This prevents the encoder from collapsing to a trivially predictable but uninformative representation.

\textbf{Formal argument.}
The total information in the latent can be decomposed as:
\begin{equation}
  I(\vz_{t+1,\text{clean}}; \vo_{t+1}) = \underbrace{I(\vz_{t+1,\text{clean}}; \vo_{t+1} \mid \mathcal{H}_t^z)}_{\text{unpredictable info (penalized)}} + \underbrace{I(\vz_{t+1,\text{clean}}; \mathcal{H}_t^z)}_{\text{history-predictable info (free)}}.
  \label{eq:info_decomp}
\end{equation}
The prior loss bounds the first term. The decoder loss incentivizes maximizing total information. The action loss ensures the latent captures dynamics-relevant structure. Together, they push the encoder to a representation that is maximally informative about action-predictable physics while discarding everything else.

\subsection{Diffusion Decoder for Future Frame Reconstruction}
\label{sec:decoder}

The decoder operates exactly as in standard UL but reconstructs the \emph{future} observation $\vo_{t+1}$. The noisy image is:
\begin{equation}
  \vo_{t+1,\tau} = \alpha_x(\tau)\,\vo_{t+1} + \sigma_x(\tau)\,\vep, \quad \vep \sim \mathcal{N}(\bm{0}, \bm{I}),
  \label{eq:forward_image}
\end{equation}
with its own log-SNR schedule $\lambda_x(\tau)$.

The decoder network $D_\theta = \hat{\vo}(\vo_{t+1,\tau}, \vz_{t+1,0}, \theta)$ conditions on both the noisy image and the slightly noisy future latent $\vz_{t+1,0}$. The reconstruction bound is:
\begin{equation}
  -\log p_\theta(\vo_{t+1} \mid \vz_{t+1,0}) \leq
  \E_{\tau \sim \mathcal{U}(0,1)}\!\left[
    \frac{d\lambda_x(\tau)}{d\tau}\,\frac{\exp\lambda_x(\tau)}{2}\,
    w_x(\lambda_x(\tau))\,
    \|\vo_{t+1} - \hat{\vo}(\vo_{t+1,\tau}, \vz_{t+1,0}, \theta)\|^2
  \right],
  \label{eq:ac_decoder_loss}
\end{equation}
with sigmoid weighting $w_x(\lambda_x(\tau)) = \sigmoid(\lambda_x(\tau) - b)$ and loss factor $c_{\text{lf}} \in [1.3, 1.7]$.

\subsection{Complete Stage 1 Objective}
\label{sec:stage1_objective}

The total per-transition loss combines the joint prior loss and the decoder loss:
\begin{equation}
  \loss_{\text{S1}}(\theta) = \underbrace{\loss_z^{\text{latent}}(\theta) + \gamma \cdot \loss_a(\theta)}_{\text{joint dynamics prior}} + c_{\text{lf}} \cdot \underbrace{\loss_x(\theta)}_{\text{decoder}}.
  \label{eq:total_loss}
\end{equation}

\textbf{Interpretation of each term:}
\begin{itemize}[leftmargin=*,itemsep=2pt]
  \item $\loss_z^{\text{latent}}$: Penalizes unpredictable information in the future latent given the history. Forces the encoder to discard action-irrelevant content.
  \item $\loss_a$: Requires the latent to contain enough dynamics structure that the action can be inferred. Prevents trivial latent collapse.
  \item $\loss_x$: Rewards information that aids reconstruction. Ensures the latent retains fine geometric detail (e.g., gripper-object distance).
\end{itemize}

\subsection{Stage 1 Training Algorithm}
\label{sec:stage1_algo}

The complete Stage~1 training procedure is given in \Cref{alg:stage1}.

\begin{algorithm}[t]
\caption{Stage 1: Training AC-UL with Joint Dynamics Prior}
\label{alg:stage1}
\begin{algorithmic}[1]
\REQUIRE Dataset of trajectories $\{(\vo_0, \va_0, \vo_1, \va_1, \ldots)\}$, history length $H$
\STATE Sample window: $(\vo_{t-H+1:t+1}, \va_{t-H:t}) \sim \mathcal{D}$
\STATE Encode all observations: $\vz_{k,\text{clean}} = E_\theta(\vo_k)$ for $k = t-H+1, \ldots, t+1$
\STATE Construct latent history: $\mathcal{H}_t^z = (\vz_{t-H+1,\text{clean}}, \va_{t-H}, \ldots, \vz_{t,\text{clean}})$
\STATE \textbf{// Joint dynamics prior loss}
\STATE Sample $\tau \sim \mathcal{U}(0,1)$, $\vep \sim \mathcal{N}(\bm{0}, \bm{I})$
\STATE $\vz_{t+1,\tau} = \alpha_z(\tau)\,\vz_{t+1,\text{clean}} + \sigma_z(\tau)\,\vep$
\STATE $(\hat{\vz}_{t+1}, \hat{\va}_t) = \hat{f}_\theta^{\text{dyn}}(\vz_{t+1,\tau}, \mathcal{H}_t^z)$
\STATE $\loss_z^{\text{latent}} = -\frac{d\lambda_z(\tau)}{d\tau}\,\frac{\exp\lambda_z(\tau)}{2}\,\|\vz_{t+1,\text{clean}} - \hat{\vz}_{t+1}\|^2 + \KL[p(\vz_{t+1,1}|\vo_{t+1})\|\mathcal{N}(\bm{0},\bm{I})]$
\STATE $\loss_a = \|\va_t - \hat{\va}_t\|^2$
\STATE \textbf{// Decoder loss}
\STATE Sample $\tau' \sim \mathcal{U}(0,1)$, $\vep' \sim \mathcal{N}(\bm{0}, \bm{I})$, $\vep_z \sim \mathcal{N}(\bm{0}, \bm{I})$
\STATE $\vz_{t+1,0} = \alpha_z(0)\,\vz_{t+1,\text{clean}} + \sigma_z(0)\,\vep_z$
\STATE $\vo_{t+1,\tau'} = \alpha_x(\tau')\,\vo_{t+1} + \sigma_x(\tau')\,\vep'$
\STATE $\loss_x = \frac{d\lambda_x(\tau')}{d\tau'}\,\frac{\exp\lambda_x(\tau')}{2}\,w_x(\lambda_x(\tau'))\,\|\vo_{t+1} - \hat{\vo}(\vo_{t+1,\tau'}, \vz_{t+1,0}, \theta)\|^2$
\STATE \textbf{// Combined loss}
\STATE $\loss_{\text{S1}} = \loss_z^{\text{latent}} + \gamma \cdot \loss_a + c_{\text{lf}} \cdot \loss_x$
\STATE Update $\theta$ by gradient descent on $\loss_{\text{S1}}$
\end{algorithmic}
\end{algorithm}

\subsection{Architecture Details}
\label{sec:architecture}

\textbf{Encoder $E_\theta$.}
A ResNet-style convolutional encoder with channel progression $[128, 256, 512, 512]$, using 2 residual blocks per downsampling stage and 3 blocks in the final stage. The encoder takes $224 \times 224$ RGB images and produces $14 \times 14 \times c$ latents with $c = 32$ channels, yielding $16\times$ spatial downsampling.

\textbf{Joint dynamics prior $\hat{f}_\theta^{\text{dyn}}$.}
A causal Vision Transformer (ViT) with 12 blocks and 768 hidden dimensions. The input sequence is constructed as an interleaved observation-action history:
\begin{equation}
  \big[\text{patchify}(\vz_{t-H+1,\text{clean}}),\; \text{mlp}(\va_{t-H}),\; \ldots,\; \text{patchify}(\vz_{t,\text{clean}}),\; \text{patchify}(\vz_{t+1,\tau})\big].
  \label{eq:prior_input}
\end{equation}
Each latent $\vz_{k,\text{clean}}$ is patchified into spatial tokens; each action $\va_k$ is projected through a 2-layer MLP into a single token. The diffusion timestep $\tau$ is encoded via sinusoidal embeddings and added to the noisy future latent tokens.

The model uses \emph{frame-level causal attention}: tokens at time $k$ can attend to all tokens at times $\leq k$ but not future times (following DreamZero). The model has two output heads:
\begin{itemize}[leftmargin=*,itemsep=2pt]
  \item \textbf{Latent head:} Linear projection from the output tokens at the $\vz_{t+1,\tau}$ positions $\to \hat{\vz}_{t+1,\text{clean}} \in \R^{h \times w \times c}$.
  \item \textbf{Action head:} MLP from the pooled output representation $\to \hat{\va}_t \in \R^{d_a}$.
\end{itemize}

\textbf{Diffusion decoder $D_\theta$.}
A UViT architecture with convolutional down/up-sampling stages $[128, 256, 512]$ and a central transformer with 8 blocks and 1024 channels. The decoder takes the noisy image $\vo_{t+1,\tau}$ and conditions on $\vz_{t+1,0}$ via adaptive layer normalization (AdaLN).

%=====================================================================
\section{Stage 2: Autoregressive Diffusion Transformer on Frozen Latents}
\label{sec:stage2}
%=====================================================================

Once Stage~1 converges, we freeze the encoder $E_\theta$ and decoder $D_\theta$. The frozen encoder provides a \emph{physics-grounded visual vocabulary} that has been optimized to contain only action-predictable geometry.

For Stage~2, we train a larger autoregressive diffusion transformer that operates entirely in the frozen latent space, following the DreamZero architecture but on AC-UL latents instead of raw video.

\subsection{Latent Trajectory Extraction}

Given a trajectory $(\vo_0, \va_0, \vo_1, \va_1, \ldots, \vo_T)$, we extract clean latents:
\begin{equation}
  \vz_{k,\text{clean}} = E_\theta(\vo_k), \quad k = 0, 1, \ldots, T.
  \label{eq:latent_extract}
\end{equation}

The trajectory in latent space becomes:
\begin{equation}
  (\vz_{0,\text{clean}}, \va_0, \vz_{1,\text{clean}}, \va_1, \ldots, \vz_{T,\text{clean}}).
  \label{eq:latent_trajectory}
\end{equation}

\subsection{Joint Video-Action Generation in Latent Space}

The Stage~2 model $\mathcal{T}_\phi$ is a diffusion transformer that takes the latent history and jointly generates the future latent and current action. We follow a chunk-wise autoregressive scheme inspired by DreamZero.

\textbf{Chunk definition.}
The trajectory is divided into chunks of $K$ latent frames. Each chunk $k$ contains:
\begin{equation}
  \mathcal{C}_k = \big(\vz_{kK+1,\text{clean}}, \va_{kK}, \vz_{kK+2,\text{clean}}, \va_{kK+1}, \ldots, \vz_{(k+1)K,\text{clean}}, \va_{(k+1)K-1}\big).
  \label{eq:chunk}
\end{equation}

\textbf{Flow-matching training.}
For each chunk, we sample noise and construct interpolants for both the latent frames and actions:
\begin{align}
  \vz_{kK+j,\tau_k} &= (1-\tau_k)\,\vz_{kK+j,\text{clean}} + \tau_k\,\vep_z^{(j)}, \label{eq:flow_latent} \\
  \va_{kK+j-1,\tau_k} &= (1-\tau_k)\,\va_{kK+j-1} + \tau_k\,\vep_a^{(j)}, \label{eq:flow_action}
\end{align}
where $\tau_k \sim \mathcal{U}(0,1)$ is sampled independently per chunk. The model predicts the velocity field:
\begin{equation}
  \loss_{\text{S2}} = \E\!\left[\frac{1}{K}\sum_{k} w(\tau_k)\,\left\|
    \mathcal{T}_\phi\!\left([\vz_{\tau_k}^k, \va_{\tau_k}^k]; \mathcal{C}_{<k}, \bm{l}, \bm{s}, \tau_k\right)
    - [\vz_{\text{clean}}^k - \vep_z^k,\; \va^k - \vep_a^k]
  \right\|^2\right],
  \label{eq:stage2_flow}
\end{equation}
where $\mathcal{C}_{<k}$ denotes the clean context from previous chunks, $\bm{l}$ is the language instruction embedding, and $\bm{s}$ is the proprioceptive state.

\textbf{Teacher forcing.}
During training, the model denoises the current chunk conditioned on \emph{clean} previous chunks via attention masking. This mirrors DreamZero's training but operates on compact AC-UL latents rather than high-dimensional video tokens.

\subsection{Closed-Loop Inference}
\label{sec:inference}

During deployment, the model operates in a closed loop:

\begin{enumerate}[leftmargin=*,itemsep=2pt]
  \item \textbf{Encode:} $\vz_{t,\text{clean}} = E_\theta(\vo_t)$ --- encode the current observation with the frozen encoder.
  \item \textbf{Generate:} $(\hat{\vz}_{t+1,\text{clean}}, \hat{\va}_t) = \mathcal{T}_\phi(\mathcal{H}_t^z)$ --- jointly generate next latent and current action via iterative denoising.
  \item \textbf{Execute:} Apply $\hat{\va}_t$ to the robot.
  \item \textbf{Observe:} Receive new observation $\vo_{t+1}$.
  \item \textbf{Replace:} Compute $\vz_{t+1,\text{clean}} = E_\theta(\vo_{t+1})$ and replace $\hat{\vz}_{t+1,\text{clean}}$ in the KV cache with the ground-truth encoding. This prevents error accumulation.
  \item \textbf{Repeat.}
\end{enumerate}

\subsection{Planning via Latent Rollouts}
\label{sec:planning}

For goal-conditioned tasks, the world model head enables planning without a separate planner. Given a goal observation $\vo_g$, we:

\begin{enumerate}[leftmargin=*,itemsep=2pt]
  \item Encode the goal: $\vz_{g,\text{clean}} = E_\theta(\vo_g)$.
  \item Use the joint model to autoregressively generate latent-action rollouts \emph{without} ground-truth replacement:
  \begin{equation}
    (\hat{\vz}_{t+1}, \hat{\va}_t) \to (\hat{\vz}_{t+2}, \hat{\va}_{t+1}) \to \cdots \to (\hat{\vz}_{t+P}, \hat{\va}_{t+P-1}).
    \label{eq:open_loop_rollout}
  \end{equation}
  \item Score each rollout by the terminal distance: $d = \|\hat{\vz}_{t+P,\text{clean}} - \vz_{g,\text{clean}}\|^2$.
  \item Use CEM or best-of-$N$ sampling to select the best action sequence.
  \item Execute the first action and replan (receding horizon).
\end{enumerate}

Because the latents encode only action-relevant geometry, the planning cost function directly measures progress toward the goal in terms of physical state change.

%=====================================================================
\section{Theoretical Analysis}
\label{sec:theory}
%=====================================================================

\subsection{AC-UL as a Variational Bound on Joint Transition Likelihood}

\begin{proposition}[AC-UL Joint Bound]
The AC-UL Stage~1 loss provides an upper bound on the negative log-likelihood of the joint observation-action transition:
\begin{equation}
  -\log p_\theta(\vo_{t+1}, \va_t \mid \mathcal{H}_t) \leq \loss_z^{\text{latent}}(\theta) + \gamma \cdot \loss_a(\theta) + c_{\text{lf}} \cdot \loss_x(\theta).
  \label{eq:acul_bound}
\end{equation}
\end{proposition}

\begin{proof}
Starting from the ELBO applied to the conditional model with latent $\vz_{t+1,0}$:
\begin{align}
  -\log p_\theta(\vo_{t+1}, \va_t \mid \mathcal{H}_t)
  &\leq \E_{\vz_{t+1,0}}\!\Big[-\log p_\theta(\vo_{t+1} \mid \vz_{t+1,0})\Big] \nonumber \\
  &\quad + \KL\!\Big[p(\vz_{t+1,0} \mid \vo_{t+1}) \;\Big\|\; p_\theta(\vz_{t+1,0} \mid \mathcal{H}_t^z)\Big] \nonumber \\
  &\quad - \E_{\vz_{t+1,0}}\!\Big[\log p_\theta(\va_t \mid \vz_{t+1,0}, \mathcal{H}_t^z)\Big].
  \label{eq:proof_step}
\end{align}
Applying the diffusion bounds (\Cref{eq:ac_prior_loss,eq:ac_decoder_loss}) to the first two terms and using $-\log p_\theta(\va_t | \cdot) \leq \gamma \|\va_t - \hat{\va}_t\|^2$ (Gaussian action likelihood) for the third term yields the result. The loss factor $c_{\text{lf}}$ loosens the bound but improves optimization by preventing posterior collapse.
\end{proof}

\subsection{Information-Theoretic Interpretation}

\textbf{Rate (prior loss).}
The prior loss $\loss_z^{\text{latent}}$ upper-bounds the mutual information between the encoder output and the data, \emph{beyond what is predictable from the history}:
\begin{equation}
  \loss_z^{\text{latent}} \geq I(\vz_{t+1,\text{clean}}; \vo_{t+1} \mid \mathcal{H}_t^z).
  \label{eq:rate_bound}
\end{equation}

\textbf{Action regularization.}
The action loss $\loss_a$ provides an additional inductive bias. The prior must recover the action from the latent transition. This means the latent must encode the \emph{causal effect} of the action on the scene:
\begin{equation}
  \loss_a \geq -I(\va_t; \vz_{t+1,\text{clean}} \mid \mathcal{H}_t^z) + \text{const},
  \label{eq:action_info}
\end{equation}
so minimizing $\loss_a$ maximizes the action information in the latent transition.

\textbf{Equilibrium.}
At convergence, the encoder achieves an optimal trade-off: it maximizes history-predictable, action-relevant information while minimizing unpredictable content:
\begin{equation}
  \vz_{t+1,\text{clean}}^* = \arg\min_{\vz} \Big[\underbrace{I(\vz; \vo_{t+1} \mid \mathcal{H}_t^z)}_{\text{rate}} + c_{\text{lf}} \cdot \underbrace{D(\vo_{t+1}, \hat{\vo}_\theta)}_{\text{distortion}} - \gamma \cdot \underbrace{I(\va_t; \vz \mid \mathcal{H}_t^z)}_{\text{action info}}\Big].
  \label{eq:rd_tradeoff}
\end{equation}

%=====================================================================
\section{Experiments}
\label{sec:experiments}
%=====================================================================

We evaluate AC-UL across three experimental settings: (1)~closed-loop planning via latent rollouts, (2)~imitation learning with joint world-action modeling, and (3)~real-world video prediction.

\subsection{Datasets and Benchmarks}
\label{sec:datasets}

\textbf{CALVIN}~\citep{mees2022calvin}.
A simulated long-horizon manipulation benchmark with 34 tasks (e.g., open drawer, push block, rotate lever). Provides RGB observations ($200 \times 200$), 7-DoF Franka Panda actions, and language instructions. We use the ABCD$\to$D split with ${\sim}600$K transitions and ${\sim}24$ hours of play data.

\textbf{LIBERO}~\citep{liu2024libero}.
A suite of 130 manipulation tasks across 5 task suites (Spatial, Object, Goal, Long, 90), featuring diverse spatial relations and long-horizon reasoning. Each task has 50 demonstrations with $128 \times 128$ RGB and 7-DoF actions.

\textbf{DROID}~\citep{khazatsky2024droid}.
A large-scale real-world dataset with ${\sim}76$K trajectories across 564 scenes on Franka Panda robots. Provides multi-view RGB ($320 \times 240$), 7-DoF actions, and language annotations. We use this for video prediction and cross-embodiment evaluation.

\textbf{Bridge-v2}~\citep{walke2023bridgedata}.
A real-world dataset with ${\sim}60$K trajectories on a WidowX robot with 7-DoF actions across kitchen tasks. Used for cross-embodiment transfer experiments.

\subsection{Baselines}
\label{sec:baselines}

\begin{itemize}[leftmargin=*,itemsep=2pt]
  \item \textbf{DreamerV3}~\citep{hafner2023dreamerv3}: Categorical latent world model with actor-critic planning.
  \item \textbf{DreamZero}~\citep{dreamzero2025}: Autoregressive diffusion transformer for joint video-action generation in pixel space. Direct comparison for the benefit of AC-UL latents vs.\ raw pixels.
  \item \textbf{DINO-WM}~\citep{zhou2024dinowm}: Frozen DINOv2 features with ViT transition model and CEM planning.
  \item \textbf{FLARE}~\citep{nvidia2025flare}: Diffusion transformer VLA with future latent alignment using SigLIP-2 features.
  \item \textbf{TD-MPC2}~\citep{hansen2024tdmpc2}: Decoder-free latent dynamics model with temporal difference learning.
  \item \textbf{IRIS}~\citep{micheli2023transformers}: Discrete autoencoder + GPT-style world model.
  \item \textbf{SD-Latent WM}: Stable Diffusion VAE encoder (frozen) + ViT dynamics model.
\end{itemize}

\subsection{Experiment 1: Closed-Loop Planning via Latent Rollouts}
\label{sec:exp_planning}

\subsubsection{Setup}

We evaluate goal-conditioned planning using the Stage~2 model's latent rollout capability (\Cref{sec:planning}). The agent receives a current observation $\vo_t$ and a goal image $\vo_g$, and must generate action sequences that reach the goal.

\textbf{Planning method.}
We use best-of-$N$ sampling: generate $N = 256$ latent-action rollouts from the joint model, score each by terminal distance to the goal latent, select the best trajectory, execute the first $K$ actions, and replan.

\textbf{Metrics.}
Success rate (SR) for discrete tasks; LPIPS between achieved and goal observations.

\subsubsection{Results}

\begin{table}[t]
\centering
\caption{Goal-conditioned planning performance (Success Rate \%). AC-UL generates action sequences via latent rollouts from the joint video-action model.}
\label{tab:planning}
\begin{tabular}{lcccccc}
\toprule
\textbf{Method} & \textbf{CALVIN-1} & \textbf{CALVIN-3} & \textbf{CALVIN-5} & \textbf{LIBERO-S} & \textbf{LIBERO-G} & \textbf{Avg.} \\
\midrule
DreamerV3          & -- & -- & -- & -- & -- & -- \\
DreamZero (pixel)  & -- & -- & -- & -- & -- & -- \\
IRIS               & -- & -- & -- & -- & -- & -- \\
DINO-WM            & -- & -- & -- & -- & -- & -- \\
TD-MPC2            & -- & -- & -- & -- & -- & -- \\
SD-Latent WM       & -- & -- & -- & -- & -- & -- \\
\midrule
AC-UL (ours)       & -- & -- & -- & -- & -- & -- \\
\bottomrule
\end{tabular}
\end{table}

\begin{table}[t]
\centering
\caption{Planning with visual generalization: novel object shapes and configurations not seen during training. Success Rate (\%).}
\label{tab:planning_generalization}
\begin{tabular}{lccc}
\toprule
\textbf{Method} & \textbf{CALVIN-Novel} & \textbf{LIBERO-Novel} & \textbf{Avg.} \\
\midrule
DreamerV3          & -- & -- & -- \\
DreamZero (pixel)  & -- & -- & -- \\
DINO-WM            & -- & -- & -- \\
AC-UL (ours)       & -- & -- & -- \\
\bottomrule
\end{tabular}
\end{table}

\subsection{Experiment 2: Imitation Learning with Joint World-Action Modeling}
\label{sec:exp_imitation}

\subsubsection{Setup}

We compare AC-UL's Stage~2 policy (the action output from the joint model) against FLARE and other VLA baselines on imitation learning benchmarks.

\textbf{Architecture.}
Stage~2 uses a 12-layer, 768-dim diffusion transformer with 12 attention heads, chunk size $K = 2$ latent frames. Language instructions are encoded with frozen CLIP. Flow-matching with $K_{\text{denoise}} = 4$ denoising steps at inference.

\textbf{Closed-loop evaluation.}
The model generates actions in closed loop: at each step, the predicted latent is replaced with the ground-truth encoding of the new observation in the KV cache.

\textbf{Metrics.}
Task success rate averaged over 3 seeds with 20 evaluation episodes per task.

\subsubsection{Results}

\begin{table}[t]
\centering
\caption{Imitation learning performance (Success Rate \%). AC-UL Stage~2 generates actions from the joint video-action model. All methods use 50 demos/task.}
\label{tab:imitation}
\begin{tabular}{lcccccc}
\toprule
\textbf{Method} & \textbf{LIBERO-S} & \textbf{LIBERO-O} & \textbf{LIBERO-G} & \textbf{LIBERO-L} & \textbf{LIBERO-90} & \textbf{Avg.} \\
\midrule
Diffusion Policy       & -- & -- & -- & -- & -- & -- \\
FLARE (policy only)    & -- & -- & -- & -- & -- & -- \\
FLARE                  & -- & -- & -- & -- & -- & -- \\
DreamZero (pixel)      & -- & -- & -- & -- & -- & -- \\
DINO-WM + BC           & -- & -- & -- & -- & -- & -- \\
\midrule
AC-UL Stage~2 (ours)   & -- & -- & -- & -- & -- & -- \\
AC-UL + planning (ours) & -- & -- & -- & -- & -- & -- \\
\bottomrule
\end{tabular}
\end{table}

\begin{table}[t]
\centering
\caption{CALVIN long-horizon evaluation. Average number of tasks completed in a chain of 5 sequential tasks. Language-conditioned.}
\label{tab:calvin}
\begin{tabular}{lccccc}
\toprule
\textbf{Method} & \textbf{1 task} & \textbf{2 tasks} & \textbf{3 tasks} & \textbf{4 tasks} & \textbf{5 tasks} \\
\midrule
HULC               & -- & -- & -- & -- & -- \\
Diffusion Policy    & -- & -- & -- & -- & -- \\
FLARE               & -- & -- & -- & -- & -- \\
DreamZero (pixel)   & -- & -- & -- & -- & -- \\
\midrule
AC-UL Stage~2 (ours) & -- & -- & -- & -- & -- \\
\bottomrule
\end{tabular}
\end{table}

\subsubsection{Data Efficiency}

\begin{table}[t]
\centering
\caption{Data efficiency on LIBERO-Spatial. Success Rate (\%) with varying demos per task.}
\label{tab:data_efficiency}
\begin{tabular}{lcccc}
\toprule
\textbf{Method} & \textbf{10 demos} & \textbf{25 demos} & \textbf{50 demos} & \textbf{100 demos} \\
\midrule
Diffusion Policy    & -- & -- & -- & -- \\
FLARE               & -- & -- & -- & -- \\
DreamZero (pixel)   & -- & -- & -- & -- \\
AC-UL Stage~2 (ours) & -- & -- & -- & -- \\
\bottomrule
\end{tabular}
\end{table}

\subsection{Experiment 3: Real-World Video-Action Prediction}
\label{sec:exp_video}

\subsubsection{Setup}

We evaluate the quality of joint video-action prediction on the DROID dataset. We train Stage~1 on DROID training split, then evaluate:
\begin{itemize}[leftmargin=*,itemsep=2pt]
  \item \textbf{Latent prediction accuracy:} MSE between predicted and ground-truth latents over multi-step rollouts (1, 5, 10, 20 steps).
  \item \textbf{Action prediction accuracy:} MSE between predicted and ground-truth actions.
  \item \textbf{Decoded video quality:} FVD, LPIPS, SSIM, PSNR of decoded predictions vs.\ ground truth.
  \item \textbf{Action-relevance metric:} Correlation between $\|\Delta \vz_k\|$ and $\|\va_k\|$.
\end{itemize}

\subsubsection{Results}

\begin{table}[t]
\centering
\caption{Video prediction quality on DROID test set. Multi-step prediction from ground-truth initial observation and open-loop rollout.}
\label{tab:video_pred}
\begin{tabular}{lcccc}
\toprule
\textbf{Method} & \textbf{FVD$\downarrow$} & \textbf{LPIPS$\downarrow$} & \textbf{SSIM$\uparrow$} & \textbf{PSNR$\uparrow$} \\
\midrule
DreamerV3            & -- & -- & -- & -- \\
DreamZero (pixel)    & -- & -- & -- & -- \\
SD-Latent WM         & -- & -- & -- & -- \\
DINO-WM              & -- & -- & -- & -- \\
\midrule
AC-UL (ours)         & -- & -- & -- & -- \\
\bottomrule
\end{tabular}
\end{table}

\begin{table}[t]
\centering
\caption{Latent prediction MSE and action prediction MSE over rollout horizon on DROID test set.}
\label{tab:latent_rollout}
\begin{tabular}{lcccc|cccc}
\toprule
& \multicolumn{4}{c|}{\textbf{Latent MSE$\downarrow$}} & \multicolumn{4}{c}{\textbf{Action MSE$\downarrow$}} \\
\textbf{Method} & \textbf{1} & \textbf{5} & \textbf{10} & \textbf{20} & \textbf{1} & \textbf{5} & \textbf{10} & \textbf{20} \\
\midrule
SD-Latent WM         & -- & -- & -- & -- & -- & -- & -- & -- \\
DINO-WM              & -- & -- & -- & -- & -- & -- & -- & -- \\
DreamZero (pixel)    & -- & -- & -- & -- & -- & -- & -- & -- \\
AC-UL (ours)         & -- & -- & -- & -- & -- & -- & -- & -- \\
\bottomrule
\end{tabular}
\end{table}

\subsection{Experiment 4: Cross-Embodiment Transfer}
\label{sec:exp_transfer}

\subsubsection{Setup}

We train Stage~1 on combined DROID (Franka) + Bridge-v2 (WidowX) data and evaluate whether a single encoder learns embodiment-agnostic physics representations.

\subsubsection{Results}

\begin{table}[t]
\centering
\caption{Cross-embodiment transfer. Stage~1 trained on combined data, Stage~2 trained on target embodiment. Success Rate (\%).}
\label{tab:cross_embodiment}
\begin{tabular}{lccc}
\toprule
\textbf{Method} & \textbf{Franka (DROID)} & \textbf{WidowX (Bridge)} & \textbf{Avg.} \\
\midrule
Separate encoders    & -- & -- & -- \\
DINO-WM (shared)     & -- & -- & -- \\
DreamZero (pixel, shared) & -- & -- & -- \\
AC-UL (shared, ours) & -- & -- & -- \\
\bottomrule
\end{tabular}
\end{table}

\subsection{Ablation Studies}
\label{sec:ablations}

\subsubsection{Effect of History-Conditioned Joint Prior}

\begin{table}[t]
\centering
\caption{Ablation: prior conditioning variants. LIBERO-Spatial planning SR (\%).}
\label{tab:ablation_prior}
\begin{tabular}{lcc}
\toprule
\textbf{Prior Variant} & \textbf{Planning SR} & \textbf{Bitrate (bpd)} \\
\midrule
Unconditional (standard UL)                      & -- & -- \\
Single-step, no action: $\vz_t \to \vz_{t+1}$   & -- & -- \\
Single-step: $(\vz_t, \va_t) \to \vz_{t+1}$     & -- & -- \\
History, latent only: $\mathcal{H}_t^z \to \vz_{t+1}$ & -- & -- \\
History, joint (AC-UL): $\mathcal{H}_t^z \to (\vz_{t+1}, \va_t)$ & -- & -- \\
\bottomrule
\end{tabular}
\end{table}

\subsubsection{Effect of Action Loss Weight}

\begin{table}[t]
\centering
\caption{Ablation: action loss weight $\gamma$ on LIBERO-Spatial.}
\label{tab:ablation_gamma}
\begin{tabular}{lccc}
\toprule
$\gamma$ & \textbf{Planning SR} & \textbf{Action MSE} & \textbf{Latent MSE} \\
\midrule
0.0 (no action loss)  & -- & -- & -- \\
0.1                    & -- & -- & -- \\
0.5                    & -- & -- & -- \\
1.0                    & -- & -- & -- \\
2.0                    & -- & -- & -- \\
\bottomrule
\end{tabular}
\end{table}

\subsubsection{Effect of History Length}

\begin{table}[t]
\centering
\caption{Ablation: history length $H$ on LIBERO-Spatial. Planning SR (\%).}
\label{tab:ablation_history}
\begin{tabular}{lccc}
\toprule
$H$ & \textbf{Planning SR} & \textbf{Action MSE} & \textbf{Latent MSE} \\
\midrule
1 (Markov)  & -- & -- & -- \\
2           & -- & -- & -- \\
4           & -- & -- & -- \\
8           & -- & -- & -- \\
\bottomrule
\end{tabular}
\end{table}

\subsubsection{Effect of Loss Factor and Noise Level}

\begin{table}[t]
\centering
\caption{Ablation: loss factor $c_{\text{lf}}$ and minimum noise level $\lambda_z(0)$ on LIBERO-Spatial.}
\label{tab:ablation_lf_noise}
\begin{tabular}{lcccc}
\toprule
$c_{\text{lf}}$ & $\lambda_z(0)$ & \textbf{Planning SR} & \textbf{PSNR} & \textbf{Bitrate} \\
\midrule
1.3 & 5   & -- & -- & -- \\
1.5 & 5   & -- & -- & -- \\
1.7 & 5   & -- & -- & -- \\
1.5 & 3   & -- & -- & -- \\
1.5 & 7   & -- & -- & -- \\
1.5 & 10  & -- & -- & -- \\
\bottomrule
\end{tabular}
\end{table}

\subsubsection{AC-UL Latents vs.\ Pixel Space for DreamZero}

\begin{table}[t]
\centering
\caption{Direct comparison: DreamZero architecture trained on AC-UL latents vs.\ raw pixels. Same model size and training budget. LIBERO success rate (\%).}
\label{tab:ablation_latent_vs_pixel}
\begin{tabular}{lccccc}
\toprule
\textbf{Input Space} & \textbf{LIBERO-S} & \textbf{LIBERO-O} & \textbf{LIBERO-G} & \textbf{Training FLOPs} & \textbf{Inference (ms)} \\
\midrule
Pixels (DreamZero)   & -- & -- & -- & -- & -- \\
SD latents           & -- & -- & -- & -- & -- \\
DINO features        & -- & -- & -- & -- & -- \\
AC-UL latents (ours) & -- & -- & -- & -- & -- \\
\bottomrule
\end{tabular}
\end{table}

%=====================================================================
\section{Analysis and Discussion}
\label{sec:analysis}
%=====================================================================

\subsection{What Do AC-UL Latents Encode?}

To verify that AC-UL latents focus on action-relevant geometry, we analyze:

\textbf{Latent sensitivity to actions.}
We compute the Jacobian $\partial \vz_{t+1,\text{clean}} / \partial \va_t$ and visualize which spatial positions are most sensitive to action perturbations. We expect high sensitivity near the gripper and manipulated objects.

\textbf{Latent change vs.\ action magnitude.}
We measure the Pearson correlation between $\|\Delta \vz_t\|$ and $\|\va_t\|$. High correlation confirms that latent dynamics are action-driven.

\textbf{Closed-loop vs.\ open-loop rollouts.}
Because AC-UL latents are compact and action-focused, we expect open-loop rollouts (without ground-truth replacement) to degrade more gracefully than pixel-space rollouts, as there is less irrelevant variation to compound.

\subsection{Comparison with DreamZero}

AC-UL and DreamZero share the autoregressive joint video-action generation paradigm but differ fundamentally in representation:
\begin{itemize}[leftmargin=*,itemsep=2pt]
  \item \textbf{DreamZero} operates on VAE-compressed video tokens. The VAE is trained for generic visual reconstruction, not action-relevant content.
  \item \textbf{AC-UL} operates on latents specifically optimized to discard action-irrelevant information. This should yield (a) more compact sequences (fewer tokens per frame), (b) better long-horizon rollout fidelity, and (c) faster inference.
\end{itemize}

\subsection{Comparison with FLARE}

FLARE adds future latent prediction as an auxiliary loss using frozen SigLIP-2 features. AC-UL differs:
\begin{itemize}[leftmargin=*,itemsep=2pt]
  \item FLARE's targets (SigLIP-2) encode general visual semantics, not action-predictable physics.
  \item AC-UL's latents are \emph{trained} to contain only action-predictable geometry.
  \item AC-UL jointly generates observations and actions; FLARE generates only actions with future features as auxiliary.
\end{itemize}

\subsection{Limitations}

\begin{itemize}[leftmargin=*,itemsep=2pt]
  \item \textbf{Two-stage training:} Stage~1 requires training three models simultaneously.
  \item \textbf{Deterministic encoder:} Inherently stochastic environments cannot be fully captured.
  \item \textbf{Fixed encoder at Stage~2:} If Stage~1 training is suboptimal, Stage~2 inherits these limitations.
  \item \textbf{Real-time inference:} The diffusion decoder adds latency for visualization; however, the action can be extracted without decoding.
\end{itemize}

%=====================================================================
\section{Related Work}
\label{sec:related}
%=====================================================================

\textbf{World models for robotics.}
Early work used recurrent state-space models~\citep{ha2018world,hafner2019dream}. DreamerV3~\citep{hafner2023dreamerv3} uses categorical latents. IRIS~\citep{micheli2023transformers} tokenizes observations with a discrete autoencoder. TD-MPC2~\citep{hansen2024tdmpc2} learns latent dynamics without a decoder. DreamZero~\citep{dreamzero2025} introduced autoregressive joint video-action generation at scale.

\textbf{Pretrained features for world models.}
DINO-WM~\citep{zhou2024dinowm} uses frozen DINOv2 patch features. R3M~\citep{nair2022r3m} and VIP~\citep{ma2023vip} learn visual representations for robotics. V-JEPA 2 uses self-supervised video features for zero-shot manipulation.

\textbf{Latent representations for generation.}
Latent Diffusion Models~\citep{rombach2022high} learn autoencoders for diffusion. Unified Latents~\citep{heek2025ul} jointly trains encoders and diffusion priors. Our work extends UL to action-conditioned dynamics.

\textbf{Vision-Language-Action models.}
RT-2~\citep{brohan2023rt2}, OpenVLA~\citep{kim2024openvla}, and GR00T N1~\citep{nvidia2025groot} fine-tune VLMs for action prediction. FLARE~\citep{nvidia2025flare} adds future latent alignment. $\pi_0$~\citep{pi02024} uses flow matching for action generation. Our Stage~2 builds on this paradigm but uses physics-grounded latents.

\textbf{Joint video-action generation.}
DreamZero~\citep{dreamzero2025} and UniSim~\citep{yang2024unisim} jointly generate video and actions. UWM~\citep{zhu2025uwm} couples video and action diffusion. Our work provides a principled latent space for such joint generation.

%=====================================================================
\section{Conclusion}
\label{sec:conclusion}
%=====================================================================

We introduced Action-Conditioned Unified Latents (AC-UL), a framework for learning visual representations optimized to encode only the physical geometry predictable from interaction history.
By extending the Unified Latents variational objective with a history-conditioned joint dynamics prior that simultaneously predicts future observations and current actions, we create an information bottleneck that discards unpredictable visual content while retaining action-relevant structure.

Our formulation follows the autoregressive video-action generation paradigm of DreamZero but operates in a principled latent space rather than raw pixels.
Stage~1 trains the physics-grounded encoder via the AC-UL objective, and Stage~2 trains an autoregressive diffusion transformer on frozen latents for joint video-action generation in closed loop.

We believe AC-UL opens promising directions: scaling Stage~2 to foundation-model sizes, incorporating action-free video for broader pretraining, and extending to multi-step chunk-wise generation for real-time control.

%=====================================================================
\bibliographystyle{plainnat}

\begin{thebibliography}{35}

\bibitem[Brohan et~al.(2023)]{brohan2023rt2}
A.~Brohan, N.~Brown, J.~Carbajal, et~al.
\newblock RT-2: Vision-language-action models transfer web knowledge to robotic control.
\newblock In \emph{CoRL}, 2023.

\bibitem[Caron et~al.(2021)]{caron2021dino}
M.~Caron, H.~Touvron, I.~Misra, et~al.
\newblock Emerging properties in self-supervised vision transformers.
\newblock In \emph{ICCV}, 2021.

\bibitem[DreamZero Team(2025)]{dreamzero2025}
DreamZero Team.
\newblock DreamZero: Autoregressive video-action world models for zero-shot robot manipulation.
\newblock \emph{arXiv preprint arXiv:2602.15922}, 2025.

\bibitem[Ha and Schmidhuber(2018)]{ha2018world}
D.~Ha and J.~Schmidhuber.
\newblock World models.
\newblock \emph{arXiv preprint arXiv:1803.10122}, 2018.

\bibitem[Hafner et~al.(2019)]{hafner2019dream}
D.~Hafner, T.~Lillicrap, I.~Fischer, et~al.
\newblock Dream to control: Learning behaviors by latent imagination.
\newblock In \emph{ICLR}, 2019.

\bibitem[Hafner et~al.(2023)]{hafner2023dreamerv3}
D.~Hafner, J.~Pasukonis, J.~Ba, and T.~Lillicrap.
\newblock Mastering diverse domains through world models.
\newblock \emph{arXiv preprint arXiv:2301.04104}, 2023.

\bibitem[Hansen et~al.(2024)]{hansen2024tdmpc2}
N.~Hansen, H.~Su, and X.~Wang.
\newblock TD-MPC2: Scalable, robust world models for continuous control.
\newblock In \emph{ICLR}, 2024.

\bibitem[Heek et~al.(2025)]{heek2025ul}
J.~Heek, E.~Hoogeboom, T.~Mensink, and T.~Salimans.
\newblock Unified Latents (UL): How to train your latents.
\newblock \emph{arXiv preprint arXiv:2602.17270}, 2025.

\bibitem[Ho et~al.(2020)]{ho2020denoising}
J.~Ho, A.~Jain, and P.~Abbeel.
\newblock Denoising diffusion probabilistic models.
\newblock In \emph{NeurIPS}, 2020.

\bibitem[Hoogeboom et~al.(2024)]{hoogeboom2024uvit}
E.~Hoogeboom, J.~Heek, and T.~Salimans.
\newblock Simple diffusion: End-to-end diffusion for high resolution images.
\newblock In \emph{ICML}, 2024.

\bibitem[James et~al.(2020)]{james2020rlbench}
S.~James, Z.~Ma, D.~Rovick~Arrojo, and A.~J.~Davison.
\newblock RLBench: The robot learning benchmark \& learning environment.
\newblock \emph{IEEE RA-L}, 2020.

\bibitem[Khazatsky et~al.(2024)]{khazatsky2024droid}
D.~Khazatsky, K.~Pertsch, S.~Nair, et~al.
\newblock DROID: A large-scale in-the-wild robot manipulation dataset.
\newblock In \emph{RSS}, 2024.

\bibitem[Kim et~al.(2024)]{kim2024openvla}
M.~J.~Kim, K.~Pertsch, S.~Karamcheti, et~al.
\newblock OpenVLA: An open-source vision-language-action model.
\newblock In \emph{CoRL}, 2024.

\bibitem[Kingma et~al.(2021)]{kingma2021variational}
D.~P.~Kingma, T.~Salimans, B.~Poole, and J.~Ho.
\newblock Variational diffusion models.
\newblock In \emph{NeurIPS}, 2021.

\bibitem[Kingma and Gao(2023)]{kingma2023understanding}
D.~P.~Kingma and R.~Gao.
\newblock Understanding diffusion objectives as the ELBO with simple data augmentation.
\newblock In \emph{NeurIPS}, 2023.

\bibitem[Liu et~al.(2024)]{liu2024libero}
B.~Liu, Y.~Zhu, C.~Gao, et~al.
\newblock LIBERO: Benchmarking knowledge transfer for lifelong robot learning.
\newblock In \emph{NeurIPS}, 2024.

\bibitem[Ma et~al.(2023)]{ma2023vip}
Y.~J.~Ma, S.~Sodhani, J.~Jayaraman, et~al.
\newblock VIP: Towards universal visual reward and representation via value-implicit pre-training.
\newblock In \emph{ICLR}, 2023.

\bibitem[Mees et~al.(2022)]{mees2022calvin}
O.~Mees, L.~Hermann, E.~Rosete-Beas, and W.~Burgard.
\newblock CALVIN: A benchmark for language-conditioned policy learning for long-horizon robot manipulation tasks.
\newblock \emph{IEEE RA-L}, 2022.

\bibitem[Micheli et~al.(2023)]{micheli2023transformers}
V.~Micheli, E.~Alonso, and F.~Fleuret.
\newblock Transformers are sample-efficient world learners.
\newblock In \emph{ICLR}, 2023.

\bibitem[Nair et~al.(2022)]{nair2022r3m}
S.~Nair, A.~Rajeswaran, V.~Kumar, et~al.
\newblock R3M: A universal visual representation for robot manipulation.
\newblock In \emph{CoRL}, 2022.

\bibitem[NVIDIA et~al.(2025a)]{nvidia2025flare}
NVIDIA, B.~Wen, S.~Yang, et~al.
\newblock FLARE: Future latent representation alignment for robot policy learning.
\newblock \emph{arXiv preprint arXiv:2505.15659}, 2025.

\bibitem[NVIDIA et~al.(2025b)]{nvidia2025groot}
NVIDIA et~al.
\newblock GR00T N1: An open foundation model for generalist humanoid robots.
\newblock \emph{arXiv preprint}, 2025.

\bibitem[Physical Intelligence(2024)]{pi02024}
Physical Intelligence.
\newblock $\pi_0$: A vision-language-action flow model for general robot control.
\newblock \emph{arXiv preprint}, 2024.

\bibitem[Rombach et~al.(2022)]{rombach2022high}
R.~Rombach, A.~Blatt, D.~Lorenz, P.~Esser, and B.~Ommer.
\newblock High-resolution image synthesis with latent diffusion models.
\newblock In \emph{CVPR}, 2022.

\bibitem[Vahdat et~al.(2021)]{vahdat2021score}
A.~Vahdat, K.~Kreis, and J.~Kautz.
\newblock Score-based generative modeling in latent space.
\newblock In \emph{NeurIPS}, 2021.

\bibitem[Walke et~al.(2023)]{walke2023bridgedata}
H.~R.~Walke, K.~Black, A.~Lee, et~al.
\newblock BridgeData V2: A dataset for robot learning at scale.
\newblock In \emph{CoRL}, 2023.

\bibitem[Yang et~al.(2024)]{yang2024unisim}
M.~Yang, Y.~Du, D.~Hafner, et~al.
\newblock Learning interactive real-world simulators.
\newblock In \emph{ICLR}, 2024.

\bibitem[Zhou et~al.(2024)]{zhou2024dinowm}
G.~Zhou, M.~Liu, et~al.
\newblock DINO-WM: World models on pre-trained visual features enable zero-shot planning.
\newblock \emph{arXiv preprint arXiv:2411.04983}, 2024.

\bibitem[Zhu et~al.(2025)]{zhu2025uwm}
C.~Zhu, Y.~Yu, et~al.
\newblock Unified World Models: Coupling video and action diffusion for pretraining on large robotic datasets.
\newblock \emph{arXiv preprint arXiv:2504.02792}, 2025.

\end{thebibliography}

%=====================================================================
\appendix
\section{Detailed Derivations}
\label{app:derivations}
%=====================================================================

\subsection{Full Derivation of the History-Conditioned Joint Prior Bound}
\label{app:prior_derivation}

\textbf{Step 1: ELBO for joint conditional generation.}
We model the joint distribution $p_\theta(\vo_{t+1}, \va_t \mid \mathcal{H}_t)$. Introducing a latent $\vz_{t+1,0}$:
\begin{align}
  -\log p_\theta(\vo_{t+1}, \va_t \mid \mathcal{H}_t)
  &= -\log \int p_\theta(\vo_{t+1} \mid \vz_{t+1,0})\, p_\theta(\va_t \mid \vz_{t+1,0}, \mathcal{H}_t^z)\, p_\theta(\vz_{t+1,0} \mid \mathcal{H}_t^z)\, d\vz_{t+1,0} \\
  &\leq -\E_{q(\vz_{t+1,0} | \vo_{t+1})}\!\left[\log p_\theta(\vo_{t+1} | \vz_{t+1,0}) + \log p_\theta(\va_t | \vz_{t+1,0}, \mathcal{H}_t^z)\right] \nonumber \\
  &\quad + \KL\!\left[q(\vz_{t+1,0} | \vo_{t+1}) \,\|\, p_\theta(\vz_{t+1,0} | \mathcal{H}_t^z)\right],
  \label{eq:app_joint_elbo}
\end{align}
where $q(\vz_{t+1,0} | \vo_{t+1}) = \mathcal{N}(\alpha_z(0)\, E_\theta(\vo_{t+1}),\, \sigma_z^2(0)\, \bm{I})$.

\textbf{Step 2: Bounding the latent KL with diffusion.}
The prior $p_\theta(\vz_{t+1,0} | \mathcal{H}_t^z)$ is a history-conditioned diffusion model. Augmenting with the full path:
\begin{align}
  &\KL\!\left[q(\vz_{t+1,0} | \vo_{t+1}) \,\|\, p_\theta(\vz_{t+1,0} | \mathcal{H}_t^z)\right] \nonumber \\
  &\leq \KL\!\left[q(\vz_{t+1,0:1} | \vo_{t+1}) \,\|\, p_\theta(\vz_{t+1,0:1} | \mathcal{H}_t^z)\right].
  \label{eq:app_chain}
\end{align}

\textbf{Step 3: Decomposing the path KL.}
Following \citet{kingma2021variational}:
\begin{align}
  &\KL\!\left[q(\vz_{t+1,0:1} | \vo_{t+1}) \,\|\, p_\theta(\vz_{t+1,0:1} | \mathcal{H}_t^z)\right] \nonumber \\
  &= \int_0^1 \E_{q(\vz_{t+1,\tau} | \vo_{t+1})}\!\left[-\frac{d\lambda_z(\tau)}{d\tau}\, \frac{\exp\lambda_z(\tau)}{2}\, \|\vz_{t+1,\text{clean}} - \hat{\vz}_\theta^{\text{dyn}}(\vz_{t+1,\tau}, \mathcal{H}_t^z)\|^2\right] d\tau \nonumber \\
  &\quad + \KL\!\left[q(\vz_{t+1,1} | \vo_{t+1}) \,\|\, \mathcal{N}(\bm{0}, \bm{I})\right].
  \label{eq:app_decompose}
\end{align}

The conditioning on $\mathcal{H}_t^z$ is the key difference from standard UL: the prior uses the full interaction history, which includes both past observation latents and past actions, to predict the future.

\textbf{Step 4: Action likelihood term.}
The action prediction component uses a Gaussian likelihood:
\begin{align}
  -\log p_\theta(\va_t | \vz_{t+1,0}, \mathcal{H}_t^z) &\propto \|\va_t - \hat{\va}_\theta(\vz_{t+1,0}, \mathcal{H}_t^z)\|^2.
  \label{eq:app_action}
\end{align}

In practice, since $\vz_{t+1,0} \approx \vz_{t+1,\text{clean}}$ (the noise at $\tau=0$ is small), we share the action prediction with the dynamics prior model evaluated at arbitrary $\tau$, which provides a training signal at all noise levels.

\textbf{Step 5: Combining all terms.}
The complete bound is:
\begin{align}
  -\log p_\theta(\vo_{t+1}, \va_t | \mathcal{H}_t) &\leq \underbrace{\loss_z^{\text{latent}}(\theta)}_{\text{latent KL bound}} + \gamma \cdot \underbrace{\loss_a(\theta)}_{\text{action MSE}} + c_{\text{lf}} \cdot \underbrace{\loss_x(\theta)}_{\text{decoder bound}}.
  \label{eq:app_total}
\end{align}

\subsection{Information-Theoretic Analysis}
\label{app:info_theory}

\textbf{Claim.}
The AC-UL prior loss upper-bounds $I(\vz_{t+1,\text{clean}}; \vo_{t+1} \mid \mathcal{H}_t^z)$.

\textbf{Proof.}
The KL in the ELBO satisfies:
\begin{align}
  \E_{p(\mathcal{H}_t)}\!\left[\KL\!\left[q(\vz_{t+1,0} | \vo_{t+1}) \,\|\, p_\theta(\vz_{t+1,0} | \mathcal{H}_t^z)\right]\right]
  &\geq I(\vz_{t+1,0}; \vo_{t+1} | \mathcal{H}_t^z).
\end{align}
When the prior is optimal ($p_\theta = p^*$), equality holds. Since $\loss_z^{\text{latent}}$ upper-bounds the KL, we have $\loss_z^{\text{latent}} \geq I(\vz_{t+1,\text{clean}}; \vo_{t+1} | \mathcal{H}_t^z)$.

\textbf{Role of joint action prediction.}
The action loss provides an additional constraint. At optimality:
\begin{equation}
  I(\va_t; \vz_{t+1,\text{clean}} | \mathcal{H}_t^z) \geq I(\va_t; \vo_{t+1} | \mathcal{H}_t^z) - I(\va_t; \vo_{t+1} | \vz_{t+1,\text{clean}}, \mathcal{H}_t^z).
\end{equation}
The action loss pushes $\vz_{t+1,\text{clean}}$ to be a \emph{sufficient statistic} for the action given the history, meaning $I(\va_t; \vo_{t+1} | \vz_{t+1,\text{clean}}, \mathcal{H}_t^z) \to 0$. This ensures the latent captures the full causal effect of the action.

%=====================================================================
\section{Implementation Details}
\label{app:implementation}
%=====================================================================

\subsection{Stage 1 Hyperparameters}

\begin{table}[h]
\centering
\caption{Stage 1 hyperparameters.}
\label{tab:app_stage1}
\begin{tabular}{ll}
\toprule
\textbf{Hyperparameter} & \textbf{Value} \\
\midrule
Image resolution & $224 \times 224$ \\
Latent spatial size & $14 \times 14$ \\
Latent channels & 32 \\
Encoder architecture & ResNet $[128, 256, 512, 512]$ \\
Prior architecture & Causal ViT-B (12 blocks, 768 dim) \\
Decoder architecture & UViT $[128, 256, 512]$ + 8-block transformer \\
History length $H$ & 4 \\
$\lambda_z(0)$ (min noise log-SNR) & 5 \\
Decoder sigmoid bias $b$ & 3.0 \\
Loss factor $c_{\text{lf}}$ & 1.5 \\
Action loss weight $\gamma$ & 1.0 \\
Optimizer & AdamW ($\beta_1 = 0.9$, $\beta_2 = 0.999$) \\
Learning rate & $1 \times 10^{-4}$ \\
LR schedule & Cosine with 5K warmup \\
Batch size & 256 \\
Training steps & 500K \\
EMA decay & 0.9999 \\
Action embedding dim & 256 \\
Action MLP & 2-layer, 256 hidden \\
Dropout & 0.1 (prior and decoder) \\
\bottomrule
\end{tabular}
\end{table}

\subsection{Stage 2 Hyperparameters}

\begin{table}[h]
\centering
\caption{Stage 2 hyperparameters.}
\label{tab:app_stage2}
\begin{tabular}{ll}
\toprule
\textbf{Hyperparameter} & \textbf{Value} \\
\midrule
Transformer layers & 12 \\
Hidden dimension & 768 \\
Attention heads & 12 \\
Chunk size $K$ & 2 latent frames \\
Context length & 8 latent frames \\
Text encoder & CLIP ViT-B/32 (frozen) \\
Training objective & Flow matching \\
Denoising steps (inference) & 4 \\
Optimizer & AdamW ($\beta_1 = 0.9$, $\beta_2 = 0.95$) \\
Learning rate & $3 \times 10^{-4}$ \\
LR schedule & Cosine with 2K warmup \\
Batch size & 256 \\
Training steps & 100K \\
Gradient clipping & 1.0 \\
\bottomrule
\end{tabular}
\end{table}

\subsection{Planning Hyperparameters}

\begin{table}[h]
\centering
\caption{Planning hyperparameters.}
\label{tab:app_planning}
\begin{tabular}{ll}
\toprule
\textbf{Hyperparameter} & \textbf{Value} \\
\midrule
Best-of-$N$ samples & 256 \\
Planning horizon $P$ & 10 steps \\
Receding horizon (execute) & $K$ = 2 steps \\
Action smoothing & Savitzky-Golay filter \\
\bottomrule
\end{tabular}
\end{table}

\end{document}
